% Kompilacja: pdflatex dokumentacja_teoretyczna.tex (dwukrotnie)
\documentclass[11pt]{article}
\usepackage[utf8]{inputenc}
\usepackage[T1]{fontenc}
\usepackage[polish]{babel}
\usepackage{lmodern}
\usepackage{amsmath,amssymb,amsthm}
\usepackage{tikz}
\usepackage[margin=2.6cm]{geometry}
\usepackage[hidelinks]{hyperref}

\theoremstyle{definition}
\newtheorem{definicja}{Definicja}[section]
\theoremstyle{plain}
\newtheorem{twierdzenie}{Twierdzenie}[section]
\newtheorem{wniosek}[twierdzenie]{Wniosek}
\theoremstyle{remark}
\newtheorem*{uwaga}{Uwaga}

% --- pomocniczy rysunek: dwa łuki nad punktami na prostej ---
% #1,#2 = końce łuku czerwonego; #3,#4 = końce łuku niebieskiego
\newcommand{\luki}[4]{%
\begin{tikzpicture}[baseline=-2pt,scale=0.52]
  \foreach \x in {1,2,3,4} { \fill (\x,0) circle (1.7pt); \node[below=2pt,font=\scriptsize] at (\x,0) {\x}; }
  \draw[red!80!black,very thick] (#1,0) to[bend left=70] (#2,0);
  \draw[blue!75!black,very thick] (#3,0) to[bend left=70] (#4,0);
\end{tikzpicture}}

\title{\large Gry kombinatoryczne\\[2pt]
\Large Arc Match --- gra Maker--Breaker na uporządkowanych skojarzeniach\\[4pt]
\large Dokumentacja teoretyczna}
\author{Wojciech Jasiński}
\date{Czerwiec 2026}

\begin{document}
\maketitle

\section{Opis gry}

\emph{Arc Match} jest dwuosobową grą typu \textbf{Maker--Breaker} rozgrywaną na
\textbf{uporządkowanym skojarzeniu}. Na prostej leży $2n$ ponumerowanych punktów. W swoim ruchu gracz
wybiera dwa jeszcze nieużyte punkty i łączy je \textbf{łukiem w swoim kolorze}. Każdy punkt może być
końcem co najwyżej jednego łuku, więc suma wszystkich narysowanych łuków jest jednym uporządkowanym
skojarzeniem, pokolorowanym wedle tego, kto narysował dany łuk.

Gracz \textbf{Maker} (czerwony) chce zbudować \emph{spośród własnych łuków} rodzinę $k$ łuków, które
\textbf{parami się przecinają} (krzyżują). Gracz \textbf{Breaker} (niebieski) stara się mu w tym
przeszkodzić. \textbf{Maker rozpoczyna grę.} Maker wygrywa w chwili, gdy jego czerwone łuki zawierają
poszukiwany wzorzec rozmiaru $k$; Breaker wygrywa, gdy plansza się zapełni (lub żadne dokończenie wzorca
przez Makera nie jest już możliwe), a wzorzec nie powstał. Remisy rozstrzygamy na korzyść Breakera ---
zgodnie ze standardową konwencją gier Maker--Breaker.

Oprócz wzorca \emph{krzyżowania} gra dopuszcza także cele alternatywne (\emph{zagnieżdżenie},
\emph{ułożenie}, dowolny pojedynczy typ) oraz tryb \emph{punktowy} (bez progu --- liczy się największa
jednorodna rodzina czerwona). W dalszej części skupiamy się na celu podstawowym, czyli krzyżowaniu,
ponieważ daje on najbardziej czytelny i strategicznie interesujący przebieg gry; dla zagnieżdżenia i
ułożenia Maker w praktyce wymusza maksymalny rozmiar niemal niezależnie od gry Breakera.

\section{Definicje}

\begin{definicja}[Uporządkowane skojarzenie]
\emph{Uporządkowanym skojarzeniem} rozmiaru $n$ nazywamy zbiór $n$ krawędzi (łuków), które parami nie
mają wspólnego końca, narysowanych na $2n$ liniowo uporządkowanych punktach $1,\dots,2n$. Każdy punkt
jest końcem dokładnie jednej krawędzi.
\end{definicja}

Definicja powyższa opisuje skojarzenie \emph{pełne} (stan końcowy, po narysowaniu wszystkich $n$
łuków). W trakcie rozgrywki plansza jest skojarzeniem \emph{częściowym} --- zbiorem dotąd narysowanych,
parami rozłącznych łuków, w którym część punktów jest jeszcze wolna.

\begin{definicja}[Trzy wzajemne położenia]
Dla dwóch krawędzi $e=(a,b)$ oraz $f=(c,d)$, gdzie $a<b$, $c<d$ i bez straty ogólności $a<c$, zachodzi
\emph{dokładnie jeden} z przypadków:
\begin{center}
\begin{tabular}{@{}lll@{}}
\textbf{Ułożenie} ($a<b<c<d$) & rozłączne, obok siebie & \luki{1}{2}{3}{4}\\[10pt]
\textbf{Zagnieżdżenie} ($a<c<d<b$) & jedno wewnątrz drugiego & \luki{1}{4}{2}{3}\\[10pt]
\textbf{Krzyżowanie} ($a<c<b<d$) & przeplatające się & \luki{1}{3}{2}{4}\\
\end{tabular}
\end{center}
\end{definicja}

\begin{definicja}[Podskojarzenie jednorodne]
Podskojarzenie jest \emph{jednorodne} typu $T$, jeśli \emph{każda} para jego krawędzi ma typ $T$.
Postaci kanoniczne rodziny jednorodnej rozmiaru $k$ (krawędzie $e_i=(a_i,b_i)$ uporządkowane wedle
$a_i$):
\[
\begin{aligned}
k\text{-krzyżowanie:}\quad & a_1<a_2<\dots<a_k< b_1<b_2<\dots<b_k,\\
k\text{-zagnieżdżenie:}\quad & a_1<a_2<\dots<a_k< b_k<\dots<b_2<b_1,\\
k\text{-ułożenie:}\quad & a_1<b_1< a_2<b_2< \dots< a_k<b_k.
\end{aligned}
\]
\end{definicja}

\paragraph{Przykład ($n=3$, punkty $1,\dots,6$).}
Skojarzenie $\{(1,4),(2,5),(3,6)\}$ daje $a\!:1<2<3 < b\!:4<5<6$, czyli \textbf{3-krzyżowanie}.
Skojarzenie $\{(1,6),(2,5),(3,4)\}$ daje $a\!:1<2<3 < b\!:6>5>4$, czyli \textbf{3-zagnieżdżenie}.
Skojarzenie $\{(1,2),(3,4),(5,6)\}$ jest \textbf{3-ułożeniem}.

\section{Twierdzenia}

Punktem wyjścia jest obserwacja, że jednorodność jest w pełnym skojarzeniu \emph{nieunikniona} ---
co czyni interesującym pytanie, czy Maker potrafi ją \emph{wymusić}, dysponując jedynie połową łuków.

\begin{twierdzenie}[Erdős--Szekeres dla skojarzeń; Dudek, Grytczuk, Ruciński~\cite{dgr-unavoidable}]
\label{tw:es}
Każde uporządkowane skojarzenie rozmiaru $n$ zawiera podskojarzenie jednorodne (krzyżowanie,
zagnieżdżenie lub ułożenie) rozmiaru co najmniej $n^{1/3}$.
\end{twierdzenie}

Grę modelujemy jako \textbf{grę pozycyjną Maker--Breaker} na hipergrafie $H$, którego wierzchołkami są
\emph{możliwe łuki}, a zbiorami wygrywającymi --- jednorodne podskojarzenia rozmiaru $k$ (Beck,
\emph{Tic-Tac-Toe Theory}~\cite{beck}). W tym języku obowiązuje klasyczne kryterium Erdősa--Selfridge'a.

\begin{twierdzenie}[Kryterium Erdősa--Selfridge'a]
\label{tw:es-selfridge}
Jeżeli dla rodziny zbiorów wygrywających $\mathcal{F}$ zachodzi
\[
\sum_{A\in\mathcal{F}} 2^{-|A|} < \tfrac12,
\]
to Breaker (gracz drugi) ma strategię wygrywającą. Ponieważ tutaj wszystkie zbiory wygrywające mają
rozmiar $k$, warunek przyjmuje postać $W(n,k,T)\cdot 2^{-k} < \tfrac12$, gdzie $W(n,k,T)$ to liczba
jednorodnych podskojarzeń rozmiaru $k$ na $2n$ punktach.
\end{twierdzenie}

\begin{uwaga}
\textbf{Istotne zastrzeżenie.} Klasyczne twierdzenie Erdősa--Selfridge'a zakłada, że elementy zbioru
bazowego są \emph{niezależnie} zajmowalne. W naszej grze zajęcie łuku \emph{zużywa jego dwa punkty},
więc łuki dzielące punkt nie mogą współistnieć --- jest to ograniczenie typu rozłączności
(matroidalne), którego podstawowe twierdzenie nie modeluje. Dlatego Twierdzenie~\ref{tw:es-selfridge}
jest tu \emph{jednostronnym warunkiem wystarczającym} i \emph{silną heurystyką} (wykorzystaną przez
sztuczną inteligencję poziomu~2), a nie gotowym rozstrzygnięciem gry. Precyzyjne uwzględnienie
rozłączności zmienia zliczanie $W(n,k,T)$ i jest właściwym kierunkiem dalszej analizy.
\end{uwaga}

\section{Próg gry $k^*(n)$}

\begin{definicja}[Próg gry]
$k^*(n)$ to największe $k$, które Maker potrafi wymusić na $2n$ punktach przy optymalnej grze obu stron.
\end{definicja}

Ponieważ Maker rysuje tylko $\lceil n/2 \rceil$ z $n$ łuków, mamy natychmiast
\[
k^*(n) \le \Big\lceil \tfrac{n}{2} \Big\rceil .
\]
Z drugiej strony Twierdzenie~\ref{tw:es} dotyczy \emph{całego} skojarzenia (obu kolorów), więc nie
ogranicza bezpośrednio $k^*(n)$ --- Maker włada jedynie połową łuków, a poszukiwany wzorzec ma być
jednobarwny. Dokładne zachowanie $k^*(n)$ pozostaje \textbf{problemem otwartym}; wartość $n^{1/3}$
traktujemy jedynie jako asymptotyczną \emph{krzywą odniesienia}. Próg $k^*(n)$ szacujemy empirycznie
symulacjami AI--vs--AI (gracze poziomu~2; dla małych $n$ można go też wyznaczyć dokładnym solverem).
Wyniki zbiera Tabela~\ref{tab:kstar}: dla $n\ge 4$ próg $k^*(n)$ leży \emph{powyżej} krzywej $n^{1/3}$ i
\emph{poniżej} pułapu $\lceil n/2\rceil$, natomiast dla $n=3$ plansza jest zbyt mała, by Maker wymusił
choćby $k=2$.

\begin{table}[h]
\centering
\begin{tabular}{c|cccccc}
$n$ & 3 & 4 & 5 & 6 & 7 & 8\\\hline
$k^*(n)$ & 0 & 2 & 2 & 2 & 3 & 3\\
$\lceil n/2\rceil$ & 2 & 2 & 3 & 3 & 4 & 4\\
$n^{1/3}$ & 1.44 & 1.59 & 1.71 & 1.82 & 1.91 & 2.00\\
\end{tabular}
\caption{Empiryczny próg $k^*(n)$ (Maker poz.~2 vs Breaker poz.~2, cel: krzyżowanie, próg wygranej
$\ge 50\%$ na 100 grach) na tle pułapu $\lceil n/2\rceil$ i krzywej nieuniknioności $n^{1/3}$.}
\label{tab:kstar}
\end{table}

\section{Podsumowanie}

Arc Match osadza klasyczną teorię gier pozycyjnych Maker--Breaker w bogatym, wizualnym kontekście
uporządkowanych skojarzeń. Nieuniknioność jednorodności (Twierdzenie~\ref{tw:es}) gwarantuje, że gra
nie jest jałowa, a kryterium Erdősa--Selfridge'a dostarcza zarówno teoretycznego ograniczenia, jak i
praktycznej heurystyki dla komputerowego przeciwnika. Najciekawszym elementem teoretycznym jest
zastrzeżenie o rozłączności łuków: jest to subtelność, której podstawowe twierdzenie nie obejmuje, oraz
naturalny punkt wyjścia do dalszych badań progu $k^*(n)$.

\begin{thebibliography}{9}
\bibitem{dgr-unavoidable}
A.~Dudek, J.~Grytczuk, A.~Ruciński,
\emph{Ordered unavoidable sub-structures in matchings and random matchings},
arXiv:2210.14042, 2022.
\bibitem{dgr-eszekeres}
A.~Dudek, J.~Grytczuk, A.~Ruciński,
\emph{Erdős--Szekeres type theorems for ordered uniform matchings},
arXiv:2301.02936, 2023.
\bibitem{dgr-twins}
A.~Dudek, J.~Grytczuk, A.~Ruciński,
\emph{Twins in ordered hyper-matchings},
arXiv:2310.01394, 2023.
\bibitem{beck}
J.~Beck,
\emph{Combinatorial Games: Tic-Tac-Toe Theory},
Cambridge University Press, 2008.
\end{thebibliography}

\end{document}
